\scanpage{191}
définie par
\[
\tau\longrightarrow\tau'=\Sigma-\tau,
\]
on aura
\[
\tau'=\operatorname{Tr}yxz=\operatorname{Tr}xzy=\operatorname{Tr}zyx.
\]
Les éléments du groupe symétrique $\mathfrak{S}_3$ échangent entre eux les neuf générateurs $\uline{b},\uline{c},\uline{t}$, et aussi $\tau$ et $\tau'$ (les permutations circ. laissant $\tau,\tau'$ fixes, les transpositions échangeant $\tau,\tau'$).
En \deleted{\unsure{termes}} tant qu'anneau où $\mathfrak{S}_3$ opère, on peut décrire $B$ comme
\[
B=k[\uline{b},\uline{b'},\uline{b''},\uline{t},\uline{t'},\uline{t''},\inserted{\uline{c},\uline{c'},\uline{c''}},\uline{\tau},\uline{\tau'}]/\bigl(\underbrace{\uline{\tau}+\uline{\tau'}-\Sigma},\underbrace{\tau\tau'-\Pi}\bigr)
\]
\margin{équations sépar\unsure{ément} invariantes sous $\mathfrak{S}_3$}
avec l'opération évidente de $\mathfrak{S}_3$ sur l'anneau des polynômes à onze variables $\uline{b},\uline{c},\uline{t},\uline{\tau},\uline{\tau'}$.

\uline{NB} Quand $2$ est inv. dans $k$, et dans ce
\margin{190}
